\section{Khảo sát thị trường và các giải pháp hiện có}\label{sec:thitruong}
Bãi đỗ xe thông minh không phải ý tưởng mới: đã có nhiều sản phẩm thương mại, cả trong
nước lẫn quốc tế. Mục này trả lời hai câu hỏi: thị trường đã có giải pháp tương tự chưa
và đã làm được đến đâu. Số liệu lấy từ báo cáo thị trường, tài liệu của nhà cung cấp và
báo chí (truy cập tháng 11/2026); thông số do nhà cung cấp tự công bố, nhóm không kiểm
chứng độc lập.

\subsection{Quy mô thị trường}
Các báo cáo đều dự báo thị trường tăng nhanh, dù con số cụ thể phụ thuộc vào phạm vi
tính của từng báo cáo. Knowledge Sourcing Intelligence ước tính thị trường bãi đỗ xe thông minh
toàn cầu đạt 8,8 tỷ USD năm 2025 và 20,4 tỷ USD năm 2030 (tăng 18,34\%/năm), trong đó
châu Á -- Thái Bình Dương tăng nhanh nhất [16]. Meticulous Research dự báo 22,9 tỷ USD vào năm
2032 (18,2\%/năm, giai đoạn 2025--2032) [17].

\subsection{Các nhóm giải pháp và mức độ hoàn thiện}
\paragraph{Phần mềm quản lý bãi giữ xe.} Đây là mức phổ biến và dễ tiếp cận nhất, với các
sản phẩm trong nước như eParking, HTG Parking, Megaparking. Chúng ghi nhận xe vào/ra,
quản lý vé và doanh thu; một số có nhận diện biển số và thanh toán trực tuyến, hoặc
dùng thẻ RFID [22]. Phần lớn chỉ quản lý lượt ra vào, không biết xe đang đỗ ở chỗ nào.

\paragraph{Kiểm soát ra vào cho tòa nhà.} Các nhà tích hợp cung cấp trọn gói thanh chắn,
đầu đọc và camera. Ví dụ, KPS System công bố hơn 800 dự án trên cả nước; hệ thống dùng
RFID, mã QR, ứng dụng di động, camera nhận diện biển số bằng AI, đầu đọc RFID tầm xa gắn
kính lái (không cần hạ kính), thanh chắn mở trong 1,8--3 giây, phân tích lưu lượng và doanh
thu, API kết nối với hệ thống quản lý tòa nhà [18]. Hikvision có giải pháp tương tự,
thêm cấu hình quy tắc tính phí, cảnh báo danh sách chặn và khớp mờ biển số [19].

\paragraph{Giám sát và dẫn hướng theo từng chỗ đỗ.} Mức này trả lời đúng câu hỏi ``chỗ nào
còn trống''. Cảm biến không dây Nedap SENSIT kết hợp từ trường và hồng ngoại để phát hiện
xe trên từng chỗ, cung cấp dữ liệu theo thời gian thực cho dẫn hướng và báo cáo [20]. SKIDATA có đèn báo từng chỗ và bảng thông tin dẫn hướng [21]; Hikvision dùng phân tích
video và cảm biến địa từ cho bãi đỗ trên phố [19].

\paragraph{Quy mô đô thị và thu phí không dừng.} Hà Nội thí điểm iParking trên điện thoại từ
2017 và đã có 33 điểm đỗ năm 2018 [23], nhưng dừng chức năng thanh toán tại 12 đơn vị từ
1/9/2020, do dự án chưa hoàn chỉnh, thiếu thiết bị kiểm soát phương tiện và nhân viên nhập
sai số xe khiến doanh thu không khớp thực tế [24]. Thu phí không dừng thì đã được áp dụng
rộng hơn: từ 5/5/2024, ô tô ra vào 5 sân bay lớn dùng tài khoản ETC [25]. Kế hoạch
198/KH-UBND ngày 22/5/2026 của Hà Nội yêu cầu từ 1/7/2026, 100\% điểm trông giữ và bãi đỗ
ô tô phải thu giá dịch vụ theo hình thức điện tử không dừng [26].

\subsection{Nhận định}
Thị trường đã có giải pháp, và ở quy mô lớn thì đã rất hoàn thiện về nhận diện xe, tốc độ
qua cổng và kết nối hệ thống. Bài học từ iParking cho thấy công nghệ chưa đủ: nếu vẫn phải
dựa vào người nhập liệu thì số liệu sẽ sai. Khoảng trống nhóm nhận thấy nằm ở bãi nhỏ:
các chức năng thường thuộc những sản phẩm riêng, phải qua nhà tích hợp và không có giá công khai
(chẳng hạn SKIDATA không niêm yết giá [27]). Đề tài hướng vào khoảng trống này, không
cạnh tranh với các hệ thống trên ở quy mô lớn.

\begin{table}[H]
\centering\small
\caption{So sánh đề tài với các nhóm giải pháp đã khảo sát}\label{tab:sosanh}
\begin{tabularx}{\linewidth}{@{}>{\raggedright\arraybackslash}p{2.7cm}XXX@{}}
\toprule
\textbf{Tiêu chí} & \textbf{Kiểm soát ra vào (KPS, Hikvision)} & \textbf{Cảm biến từng chỗ (Nedap, SKIDATA)} & \textbf{Đề tài} \\ \midrule
Nhận diện xe & Biển số, RFID tầm xa, QR & Không (chỉ biết có xe) & RFID tầm gần, không có biển số \\
Biết từng chỗ trống & Không thấy trong tài liệu của KPS & Có, là chức năng chính & Có, 5 cảm biến hồng ngoại \\
Tính phí, thanh toán & Có, theo nhiều quy tắc & Cần hệ thống khác & Tự tính phí; VietQR tự xác nhận qua SePay \\
Cảnh báo cháy, ngập gắn với thanh chắn & Không thấy trong tài liệu đã khảo sát & Không thấy trong tài liệu đã khảo sát & Có: cháy thì mở cả hai thanh chắn \\
Máy chủ, phần mềm & Máy chủ và phần mềm quản lý & Cần nền tảng quản lý & Chạy ngay trên ESP32-C3 \\
Quy mô, độ chín & Dự án thực tế, nhiều năm vận hành & Triển khai ở quy mô đô thị & Mô hình 5 chỗ, 2 cổng \\
Chi phí & Báo giá theo dự án & Báo giá theo dự án & Linh kiện phổ thông, nhóm tự lắp \\
\bottomrule
\end{tabularx}
\end{table}

\section{Lợi thế của giải pháp}\label{sec:loithe}
So với các nhóm trên, giải pháp của nhóm có những điểm mạnh sau:
\begin{enumerate}
  \item \textbf{Gom đủ chức năng trong một hệ thống.} Giám sát từng chỗ, kiểm soát ra vào,
        tính phí, thanh toán và cảnh báo an toàn chạy chung trên ba bo mạch, không phải ghép
        nhiều sản phẩm.
  \item \textbf{An toàn gắn trực tiếp với thanh chắn.} Khi phát hiện lửa, hệ thống tự bật
        còi, hiển thị cảnh báo và mở cả hai thanh chắn để sơ tán, không chờ người vận hành.
        Đây là điểm nhóm không thấy nêu trong tài liệu sản phẩm đã khảo sát.
  \item \textbf{Thanh toán phù hợp người dùng Việt Nam.} Người lái quét VietQR bằng ứng dụng
        ngân hàng hoặc ví quen thuộc. Bãi xe không cần cổng thanh toán riêng, người lái không
        cần cài thêm ứng dụng; thanh chắn ra chỉ mở khi đã nhận tiền. Chưa thanh toán thì
        thanh chắn vẫn đóng, và người vận hành có thể xác nhận thủ công.
  \item \textbf{Giảm phụ thuộc vào người nhập liệu.} Thời gian và phí do hệ thống tự tính,
        giao dịch được đối chiếu tự động, tránh đúng loại sai lệch doanh thu đã gặp ở iParking.
  \item \textbf{Không cần máy chủ riêng hay cài ứng dụng.} Toàn bộ logic chạy trên ESP32-C3;
        trang giám sát mở bằng trình duyệt tại \code{http://parking.local}.
  \item \textbf{Vẫn chạy khi mất Internet.} Chỉ riêng bước tự xác nhận chuyển khoản phải
        chuyển sang người vận hành; các chức năng khác không bị ảnh hưởng.
  \item \textbf{Chi phí thấp, dễ thay thế.} Dùng linh kiện phổ thông như ESP8266, RC522,
        SG90 và cảm biến hồng ngoại, nên hỏng ở đâu thay ở đó.
  \item \textbf{Kiểm chứng được.} 397 phép kiểm thử tự động chạy trên máy tính, không cần
        phần cứng (Chương~\ref{chap:kiemthu}).
\end{enumerate}

\paragraph{Những điểm chưa bằng hệ thống thương mại.} Đề tài chưa có nhận diện biển số và đầu
đọc tầm xa nên xe phải dừng để quẹt thẻ; chưa có ETC nên chưa đáp ứng yêu cầu thu không
dừng của Hà Nội đối với bãi đỗ ô tô [26]; xác nhận tự động phụ thuộc vào dịch vụ SePay; hệ
thống mới được thử trên mô hình, chưa qua vận hành thực tế. Các hạn chế còn lại được nêu ở
Mục~\ref{sec:conlai}.
